\begin{problem}[Challenge: chains in a local spectrum]
Let $\mathfrak p\in\Spec A$.  Prove that chains of prime ideals in $A_{\mathfrak p}$ correspond bijectively to chains
\[
\mathfrak q_0\subsetneq\cdots\subsetneq\mathfrak q_n\subseteq\mathfrak p
\]
in $A$.  Explain how this anticipates the formula $\dim A_{\mathfrak p}=\operatorname{ht}(\mathfrak p)$.
\end{problem}

\paragraph{Why this problem matters.}
It connects local spectra with the later dimension theory while staying at the level of prime chains.

\paragraph{Useful ideas before starting.}
Use extension-contraction under localization and the inclusion-preserving prime correspondence.

\paragraph{Guided hints.}
Every prime of $A_{\mathfrak p}$ contracts to a prime below $\mathfrak p$, and strictness is preserved.

\begin{solution}
The localization correspondence is an inclusion-preserving bijection between $\Spec A_{\mathfrak p}$ and primes $\mathfrak q\subseteq\mathfrak p$.  Therefore it carries strict chains to strict chains of the same length in both directions.  Taking suprema of lengths gives precisely the statement that the dimension of the local ring equals the supremum of chain lengths ending at $\mathfrak p$, which is the height of $\mathfrak p$.
\end{solution}

\paragraph{What happened mathematically.}
Local dimension counts how many specialization steps can occur before reaching the chosen point.

\paragraph{Extensions.}
Apply the result to $k[x_1,\dots,x_n]$ at the origin.

\paragraph{Provenance.}
New canonical synthesis problem based on the local-ring and residue-field corpus.
